\documentclass[12pt]{article}
\usepackage{amsmath, amssymb}
\usepackage[margin=1in]{geometry}
\setlength{\parskip}{6pt}
\title{QUBO Formulation: Balanced Distance-Based Clustering Problem}
\date{}
\begin{document}
\maketitle

\subsection*{Balanced Distance-Based Clustering Problem}

\paragraph{Problem Statement.}
Given $N$ data points and a target number of clusters $K$, the task is to assign each point to exactly one cluster such that every cluster contains exactly $N/K$ points. The objective is to minimize the total within-cluster distance, defined as the sum of pairwise distances $d_{i,j}$ for all pairs of points $\{i,j\}$ that are assigned to the same cluster.

\paragraph{Decision Variables.}
We introduce binary decision variables $x_{i,c} \in \{0,1\}$ for each point $i \in \{0, \dots, N-1\}$ and each cluster $c \in \{0, \dots, K-1\}$. The variable $x_{i,c}$ equals $1$ if point $i$ is assigned to cluster $c$, and $0$ otherwise.

\paragraph{QUBO Derivation.}
Step 1 --- Original Objective. The problem is naturally a minimization problem:
\begin{align}
\min \quad \sum_{0 \le i < j \le N-1} d_{i,j} \sum_{c=0}^{K-1} x_{i,c} x_{j,c}.
\end{align}
Since QUBO solvers minimize a quadratic form $H(x) = x^T Q x + c_0$, no sign flip is required.

Step 2 --- Constraints. The assignment and cluster-size requirements are expressed as:
\begin{align}
\sum_{c=0}^{K-1} x_{i,c} &= 1, \quad \forall i \in \{0, \dots, N-1\}, \\
\sum_{i=0}^{N-1} x_{i,c} &= \frac{N}{K}, \quad \forall c \in \{0, \dots, K-1\}.
\end{align}

Step 3 --- Penalty Conversion. Let $\lambda > 0$ be a penalty weight. We convert each constraint into a quadratic penalty term by squaring the constraint residual.

For the assignment constraint, the penalty contribution for point $i$ is:
\begin{align}
P_{\text{assign},i} &= \lambda \left( \sum_{c=0}^{K-1} x_{i,c} - 1 \right)^2 \\
&= \lambda \left( \sum_{c=0}^{K-1} x_{i,c}^2 + \sum_{c \neq c'} x_{i,c} x_{i,c'} - 2 \sum_{c=0}^{K-1} x_{i,c} + 1 \right).
\end{align}
Using the idempotent property $x_{i,c}^2 = x_{i,c}$, the quadratic and linear terms combine as $\sum_{c=0}^{K-1} x_{i,c}^2 - 2 \sum_{c=0}^{K-1} x_{i,c} = -\sum_{c=0}^{K-1} x_{i,c}$. Thus:
\begin{align}
P_{\text{assign},i} = \lambda \left( \sum_{c \neq c'} x_{i,c} x_{i,c'} - \sum_{c=0}^{K-1} x_{i,c} + 1 \right).
\end{align}

For the cluster-size constraint, the penalty contribution for cluster $c$ is:
\begin{align}
P_{\text{size},c} &= \lambda \left( \sum_{i=0}^{N-1} x_{i,c} - \frac{N}{K} \right)^2 \\
&= \lambda \left( \sum_{i=0}^{N-1} x_{i,c}^2 + \sum_{i \neq i'} x_{i,c} x_{i',c} - 2\frac{N}{K} \sum_{i=0}^{N-1} x_{i,c} + \left(\frac{N}{K}\right)^2 \right).
\end{align}
Applying $x_{i,c}^2 = x_{i,c}$ yields $\sum_{i=0}^{N-1} x_{i,c}^2 - 2\frac{N}{K} \sum_{i=0}^{N-1} x_{i,c} = -\left(2\frac{N}{K} - 1\right) \sum_{i=0}^{N-1} x_{i,c}$. Thus:
\begin{align}
P_{\text{size},c} = \lambda \left( \sum_{i \neq i'} x_{i,c} x_{i',c} - \left(2\frac{N}{K} - 1\right) \sum_{i=0}^{N-1} x_{i,c} + \left(\frac{N}{K}\right)^2 \right).
\end{align}

Step 4 --- Combined Hamiltonian. The full QUBO energy function $H(x)$ is the sum of the objective and all penalty terms:
\begin{align}
H(x) = \sum_{0 \le i < j \le N-1} d_{i,j} \sum_{c=0}^{K-1} x_{i,c} x_{j,c} + \sum_{i=0}^{N-1} P_{\text{assign},i} + \sum_{c=0}^{K-1} P_{\text{size},c}.
\end{align}

Step 5 --- Q Matrix and Offset. Mapping variables to linear indices $k = iK + c$, we read off the coefficients of $x_k$ and $x_k x_l$ to obtain the QUBO parameters:
\begin{align}
Q_{(i,c),(j,c)} &= d_{i,j}, \quad \forall 0 \le i < j \le N-1, \\
Q_{(i,c),(i,c')} &= \lambda, \quad \forall i \in \{0,\dots,N-1\}, \; c \neq c', \\
Q_{(i,c),(i',c)} &= \lambda, \quad \forall i \neq i', \; c \in \{0,\dots,K-1\}, \\
Q_{(i,c),(i,c)} &= -\lambda - \lambda\left(2\frac{N}{K} - 1\right) = -2\lambda \frac{N}{K}, \\
c_0 &= \lambda \sum_{i=0}^{N-1} 1 + \lambda \sum_{c=0}^{K-1} \left(\frac{N}{K}\right)^2 = \lambda N \left(1 + \frac{N}{K}\right).
\end{align}
All unspecified off-diagonal entries are zero.

\paragraph{Penalty Weight Justification.}
The penalty weight $\lambda$ must be chosen strictly greater than the maximum possible reduction in the objective function that could result from violating a constraint. Since distances $d_{i,j}$ are non-negative and the objective is a sum of at most $\binom{N}{2}$ terms, the maximum feasible objective value is bounded by $D_{\max} \binom{N}{2}$, where $D_{\max} = \max_{i,j} d_{i,j}$. Choosing $\lambda > D_{\max} \binom{N}{2}$ guarantees that any constraint violation incurs a penalty strictly larger than any feasible objective improvement, thereby enforcing feasibility at the global minimum.

\paragraph{Correctness Argument.}
Minimizing $H(x)$ over $\{0,1\}^{NK}$ recovers the optimal clustering because $H(x)$ equals the original objective function for all feasible assignments (where all penalty terms vanish), and strictly exceeds the minimum feasible objective for any infeasible assignment (due to the positive penalty terms). Specifically, if $x$ violates any constraint, $H(x) \ge \lambda$. If $\lambda$ is chosen as justified above, then $H(x) > \min_{x \in \text{Feasible}} \text{Objective}(x)$. Consequently, the global minimizer of $H(x)$ must satisfy all constraints and simultaneously minimize the original objective, yielding the optimal balanced clustering.

\paragraph{Illustrative Example.}
Consider $N=4$ points and $K=2$ clusters, so each cluster must contain exactly $2$ points. The variables are $x_{i,c}$ for $i \in \{0,1,2,3\}$ and $c \in \{0,1\}$, mapped to linear indices $k = 2i + c$. Let $\lambda = 100$. The concrete objective terms are:
\begin{align}
\text{Obj}(x) &= d_{01}(x_{0,0}x_{1,0} + x_{0,1}x_{1,1}) + d_{02}(x_{0,0}x_{2,0} + x_{0,1}x_{2,1}) + d_{03}(x_{0,0}x_{3,0} + x_{0,1}x_{3,1}) \\
&\quad + d_{12}(x_{1,0}x_{2,0} + x_{1,1}x_{2,1}) + d_{13}(x_{1,0}x_{3,0} + x_{1,1}x_{3,1}) + d_{23}(x_{2,0}x_{3,0} + x_{2,1}x_{3,1}).
\end{align}
The assignment penalties for each point $i$ expand to $\lambda \left( x_{i,0}x_{i,1} - x_{i,0} - x_{i,1} + 1 \right)$, yielding linear coefficients of $-100$ and a cross term of $100$ for each $i$. The cluster-size penalties for each cluster $c$ expand to $\lambda \left( \sum_{i \neq i'} x_{i,c}x_{i',c} - 3\sum_{i=0}^{3} x_{i,c} + 4 \right)$, yielding linear coefficients of $-300$ and cross terms of $100$ for each $c$. Combining these with the objective gives the full Hamiltonian $H(x)$. The diagonal entries are $Q_{k,k} = -2(100)(4/2) = -400$ for all $k \in \{0,\dots,7\}$. The offset is $c_0 = 100 \cdot 4 (1 + 2) = 1200$. The resulting QUBO matrix entries are explicitly:
\begin{align}
Q_{k,k} &= -400, \quad \forall k \in \{0,\dots,7\}, \\
Q_{(i,c),(j,c)} &= d_{i,j}, \quad \forall 0 \le i < j \le 3, \\
Q_{(i,c),(i,c')} &= 100, \quad \forall i \in \{0,\dots,3\}, c \neq c', \\
Q_{(i,c),(i',c)} &= 100, \quad \forall i \neq i', c \in \{0,1\},
\end{align}
with all other off-diagonal entries equal to $0$. For a concrete distance matrix, the optimal assignment minimizes the sum of intra-cluster distances while satisfying the balance constraint. For instance, if points are arranged linearly with distances proportional to their indices, the optimal clustering groups $\{0,1\}$ and $\{2,3\}$, yielding a minimum energy of $H(x^*) = d_{01} + d_{23} + 1200$.

\end{document}